\textbf{Phần 1. Trắc nghiệm lựa chọn}
\setcounterpageref{ex}{0}

\begin{ex} %[5D10N60]
Công thức tính quãng đường $s$ khi biết vận tốc $v$ và thời gian $t$ là
\choice
{$s = v : t$}
{$s = t : v$}
{$s = v + t$}
{\True $s = v \times t$}
\loigiai{
$s = v \times t$.\\
\textbf{Đáp án của câu hỏi là D}
}
\end{ex}

\begin{ex} %[5D10N60]
Công thức tính thời gian $t$ khi biết quãng đường $s$ và vận tốc $v$ là
\choice
{$t = s \times v$}
{\True $t = s : v$}
{$t = v : s$}
{$t = s - v$}
\loigiai{
$t = s : v$.\\
\textbf{Đáp án của câu hỏi là B}
}
\end{ex}

\begin{ex} %[5D10H60]
Một người đi xe máy với vận tốc $38\text{ km/h}$ trong $3$ giờ. Quãng đường người đó đi được là
\choice
{$104\text{ km}$}
{$124\text{ km}$}
{\True $114\text{ km}$}
{$110\text{ km}$}
\loigiai{
$s = 38 \times 3 = 114\text{ (km)}$.\\
\textbf{Đáp án của câu hỏi là C}
}
\end{ex}

\begin{ex} %[5D10H60]
Một ô tô đi quãng đường $150\text{ km}$ với vận tốc $50\text{ km/h}$. Thời gian ô tô đi hết quãng đường là
\choice
{\True $3\text{ giờ}$}
{$2\text{ giờ}$}
{$4\text{ giờ}$}
{$2{,}5\text{ giờ}$}
\loigiai{
$t = 150 : 50 = 3\text{ (giờ)}$.\\
\textbf{Đáp án của câu hỏi là A}
}
\end{ex}

\begin{ex} %[5D10H60]
Một người đi bộ với vận tốc $4{,}5\text{ km/h}$ trong $1$ giờ $20$ phút. Quãng đường người đó đi được là
\choice
{$5\text{ km}$}
{$6{,}5\text{ km}$}
{$7\text{ km}$}
{\True $6\text{ km}$}
\loigiai{
$1$ giờ $20$ phút $= \dfrac{4}{3}$ giờ.\\
$s = 4{,}5 \times \dfrac{4}{3} = 6\text{ (km)}$.\\
\textbf{Đáp án của câu hỏi là D}
}
\end{ex}

\begin{ex} %[5D10H60]
Một người đi xe đạp quãng đường $27\text{ km}$ với vận tốc $12\text{ km/h}$. Thời gian người đó đi là
\choice
{$2$ giờ $10$ phút}
{$2$ giờ $20$ phút}
{\True $2$ giờ $15$ phút}
{$2$ giờ $30$ phút}
\loigiai{
$t = 27 : 12 = 2{,}25\text{ giờ} = 2$ giờ $15$ phút.\\
\textbf{Đáp án của câu hỏi là C}
}
\end{ex}

\begin{ex} %[5D10V60]
Hai thành phố A và B cách nhau $160\text{ km}$. Hai ô tô khởi hành cùng một lúc từ A và B đi ngược chiều nhau. Xe thứ nhất đi với vận tốc $45\text{ km/h}$, xe thứ hai đi với vận tốc $35\text{ km/h}$. Sau bao lâu hai xe gặp nhau?
\choice
{$1{,}5\text{ giờ}$}
{$2{,}5\text{ giờ}$}
{$3\text{ giờ}$}
{\True $2\text{ giờ}$}
\loigiai{
Tổng vận tốc hai xe: $45 + 35 = 80\text{ (km/h)}$.\\
Thời gian gặp nhau: $160 : 80 = 2\text{ (giờ)}$.\\
\textbf{Đáp án của câu hỏi là D}
}
\end{ex}

\begin{ex} %[5D10V60]
Một người đi xe máy khởi hành từ A lúc $7$ giờ với vận tốc $36\text{ km/h}$. Đến $8$ giờ, một ô tô cũng xuất phát từ A đuổi theo xe máy với vận tốc $54\text{ km/h}$. Ô tô đuổi kịp xe máy lúc mấy giờ?
\choice
{$9$ giờ}
{\True $10$ giờ}
{$10$ giờ $30$ phút}
{$11$ giờ}
\loigiai{
Trong 1 giờ (từ 7h đến 8h), xe máy đi được: $36 \times 1 = 36\text{ (km)}$.\\
Hiệu vận tốc hai xe: $54 - 36 = 18\text{ (km/h)}$.\\
Thời gian ô tô đuổi kịp: $36 : 18 = 2\text{ (giờ)}$.\\
Thời điểm gặp nhau: $8 + 2 = 10\text{ (giờ)}$.\\
\textbf{Đáp án của câu hỏi là B}
}
\end{ex}

\begin{ex} %[5D10V60]
Một ca nô đi xuôi dòng từ bến A đến bến B hết $2$ giờ, đi ngược dòng từ B về A hết $3$ giờ. Biết vận tốc dòng nước là $3\text{ km/h}$. Quãng đường sông từ A đến B là
\choice
{\True $36\text{ km}$}
{$30\text{ km}$}
{$40\text{ km}$}
{$45\text{ km}$}
\loigiai{
Hiệu vận tốc xuôi dòng và ngược dòng bằng 2 lần vận tốc dòng nước: $3 \times 2 = 6\text{ (km/h)}$.\\
Tỉ số thời gian xuôi và ngược là $\dfrac{2}{3} \Rightarrow$ Tỉ số vận tốc xuôi và ngược là $\dfrac{3}{2}$.\\
Vận tốc xuôi dòng: $6 : (3 - 2) \times 3 = 18\text{ (km/h)}$.\\
Quãng đường AB: $18 \times 2 = 36\text{ (km)}$.\\
\textbf{Đáp án của câu hỏi là A}
}
\end{ex}

\begin{ex} %[5D10V60]
Một người đi xe đạp từ A đến B với vận tốc $12\text{ km/h}$. Sau đó đi ô tô từ B về A với vận tốc $48\text{ km/h}$. Tổng thời gian cả đi và về là $2$ giờ $30$ phút. Quãng đường AB là
\choice
{$20\text{ km}$}
{$28\text{ km}$}
{\True $24\text{ km}$}
{$30\text{ km}$}
\loigiai{
Tỉ số vận tốc đi và về: $12 : 48 = \dfrac{1}{4} \Rightarrow$ Tỉ số thời gian đi và về là $\dfrac{4}{1}$.\\
Đổi $2$ giờ $30$ phút $= 2{,}5$ giờ.\\
Thời gian đi xe đạp: $2{,}5 : (4 + 1) \times 4 = 2$ giờ.\\
Quãng đường AB: $12 \times 2 = 24\text{ (km)}$.\\
\textbf{Đáp án của câu hỏi là C}
}
\end{ex}

\begin{ex} %[5D10V60]
Một đoàn tàu hỏa dài $150\text{ m}$ chạy qua một chiếc cầu dài $450\text{ m}$ với vận tốc $36\text{ km/h}$. Thời gian để đoàn tàu chạy hoàn toàn qua cây cầu là
\choice
{\True $60\text{ giây}$}
{$45\text{ giây}$}
{$50\text{ giây}$}
{$75\text{ giây}$}
\loigiai{
Quãng đường tàu phải chạy: $450 + 150 = 600\text{ (m)}$.\\
Vận tốc $36\text{ km/h} = 10\text{ m/s}$.\\
Thời gian tàu qua cầu: $600 : 10 = 60\text{ (giây)}$.\\
\textbf{Đáp án của câu hỏi là A}
}
\end{ex}

\begin{ex} %[5D10C60]
Hai người khởi hành cùng lúc từ hai địa điểm A và B cách nhau $72\text{ km}$ và đi ngược chiều nhau. Người thứ nhất đi từ A bằng xe đạp với vận tốc $14\text{ km/h}$, người thứ hai đi từ B bằng xe máy với vận tốc $34\text{ km/h}$. Hỏi sau bao lâu thì hai người cách nhau $24\text{ km}$ trước khi gặp nhau?
\choice
{$45\text{ phút}$}
{\True $1\text{ giờ}$}
{$1{,}2\text{ giờ}$}
{$1{,}5\text{ giờ}$}
\loigiai{
Tổng quãng đường hai người đã đi được: $72 - 24 = 48\text{ (km)}$.\\
Tổng vận tốc: $14 + 34 = 48\text{ (km/h)}$.\\
Thời gian: $48 : 48 = 1\text{ (giờ)}$.\\
\textbf{Đáp án của câu hỏi là B}
}
\end{ex}

\textbf{Phần 2. Trắc nghiệm đúng sai}
\setcounterpageref{ex}{0}

\begin{ex} %[5D10H60]
Quãng đường AB dài $120\text{ km}$. Một ô tô đi từ A đến B với vận tốc $50\text{ km/h}$. Xét tính đúng sai:
\choiceTF
{\True Thời gian ô tô đi hết quãng đường là $2{,}4$ giờ.}
{Đổi $2{,}4$ giờ $= 2$ giờ $40$ phút.}
{\True Đổi $2{,}4$ giờ $= 2$ giờ $24$ phút.}
{\True Nếu xuất phát lúc $8$ giờ thì ô tô đến B lúc $10$ giờ $24$ phút.}
\loigiai{
\begin{itemize}
    \item $t = 120 : 50 = 2{,}4$ giờ. Đúng.
    \item $2{,}4$ giờ $= 2$ giờ $+ 0{,}4 \times 60$ phút $= 2$ giờ $24$ phút $\neq 2$ giờ $40$ phút. Ý 2 Sai, Ý 3 Đúng.
    \item Đến nơi lúc $8$h $+ 2$h $24$p $= 10$h $24$p. Đúng.
\end{itemize}
\begin{itemize}
    \item a) \textbf{Đúng}.
    \item b) \textbf{Sai}.
    \item c) \textbf{Đúng}.
    \item d) \textbf{Đúng}.
\end{itemize}
}
\end{ex}

\begin{ex} %[5D10V60]
Hai thành phố A và B cách nhau $135\text{ km}$. Một xe máy đi từ A đến B với vận tốc $30\text{ km/h}$, cùng lúc đó một ô tô đi từ B về A với vận tốc $60\text{ km/h}$. Xét tính đúng sai:
\choiceTF
{\True Sau mỗi giờ, hai xe tiến lại gần nhau thêm $90\text{ km}$.}
{\True Hai xe gặp nhau sau $1{,}5$ giờ ($1$ giờ $30$ phút).}
{Nơi hai xe gặp nhau cách A là $60\text{ km}$.}
{\True Nơi hai xe gặp nhau cách A là $45\text{ km}$.}
\loigiai{
\begin{itemize}
    \item Tổng vận tốc: $30 + 60 = 90\text{ km/h}$. Đúng.
    \item Thời gian gặp: $135 : 90 = 1{,}5$ giờ. Đúng.
    \item Điểm gặp cách A: $30 \times 1{,}5 = 45\text{ km} \neq 60\text{ km}$. Ý 3 Sai, Ý 4 Đúng.
\end{itemize}
\begin{itemize}
    \item a) \textbf{Đúng}.
    \item b) \textbf{Đúng}.
    \item c) \textbf{Sai}.
    \item d) \textbf{Đúng}.
\end{itemize}
}
\end{ex}

\textbf{Phần 3. Trả lời ngắn}
\setcounterpageref{ex}{0}

\begin{ex} %[5D10H60]
Một người đi xe máy với vận tốc $42\text{ km/h}$ trong $2{,}5$ giờ. Tính quãng đường đi được theo ki-lô-mét.
\shortans{105}
\loigiai{
\textbf{Đáp số:} $105$\\
$42 \times 2{,}5 = 105\text{ km}$.
}
\end{ex}

\begin{ex} %[5D10H60]
Một ô tô đi quãng đường $168\text{ km}$ với vận tốc $48\text{ km/h}$. Tính thời gian ô tô đi theo giờ (dùng dấu phẩy nếu có).
\shortans{3,5}
\loigiai{
\textbf{Đáp số:} $3{,}5$\\
$168 : 48 = 3{,}5\text{ giờ}$.
}
\end{ex}

\begin{ex} %[5D10V60]
Một người đi bộ với vận tốc $5\text{ km/h}$ từ nhà đến trường hết $24$ phút. Tính quãng đường từ nhà đến trường theo ki-lô-mét.
\shortans{2}
\loigiai{
\textbf{Đáp số:} $2$\\
$24$ phút $= 0{,}4$ giờ $\Rightarrow 5 \times 0{,}4 = 2\text{ km}$.
}
\end{ex}

\begin{ex} %[5D10V60]
Hai ô tô khởi hành cùng lúc từ hai điểm cách nhau $180\text{ km}$ đi ngược chiều nhau với vận tốc $48\text{ km/h}$ và $42\text{ km/h}$. Hỏi sau bao nhiêu giờ hai xe gặp nhau?
\shortans{2}
\loigiai{
\textbf{Đáp số:} $2$\\
$180 : (48 + 42) = 180 : 90 = 2\text{ giờ}$.
}
\end{ex}

\begin{ex} %[5D10V60]
Một ca nô đi xuôi dòng một khúc sông với vận tốc $24\text{ km/h}$ hết $2$ giờ $15$ phút. Tính độ dài khúc sông đó theo ki-lô-mét.
\shortans{54}
\loigiai{
\textbf{Đáp số:} $54$\\
$24 \times 2{,}25 = 54\text{ km}$.
}
\end{ex}

\begin{ex} %[5D10C60]
Một người đi xe máy từ A đến B với vận tốc $35\text{ km/h}$. Sau $1$ giờ, một ô tô cũng đi từ A đuổi theo với vận tốc $50\text{ km/h}$. Hỏi sau bao nhiêu phút ô tô đuổi kịp xe máy?
\shortans{140}
\loigiai{
\textbf{Đáp số:} $140$\\
Khoảng cách ban đầu: $35\text{ km}$.\\
Thời gian đuổi kịp: $35 : (50 - 35) = \dfrac{35}{15} = \dfrac{7}{3}\text{ giờ} = 140\text{ phút}$.
}
\end{ex}
